\documentclass[AER]{AEA}

% --- 必须的宏包 ---
\usepackage{amsmath, amssymb, mathtools}
\usepackage{newtxtext}
\usepackage{newtxmath}
\usepackage{bm}
\usepackage{booktabs}
\usepackage{array}
\usepackage{placeins} 
\usepackage{natbib} 
\usepackage{url}
\usepackage{xcolor} % 确保颜色宏包已加载
\usepackage[hidelinks]{hyperref}
\usepackage{tikz}

% 【修复】：给 blacksquare 加上数学模式符号 $ $
\newcommand{\qed}{\nolinebreak\hfill\ensuremath{\blacksquare}\par}

\newcommand{\noopsort}[1]{} 
% 【关键修改】：在导言区定义为空，这意味着在正文中调用时，它什么也不显示！
\DeclareRobustCommand{\AERlink}[1]{}

\urlstyle{same}

% --- 核心设置 ---
\issueName{THE AMERICAN ECONOMIC RUBBISH}
\draftSpacing{1.5}

\begin{document}

% --- 封面信息 ---
% 1. 不使用 \thanks，在标题末尾直接硬编码一个上标十字架
\title{Why Working is a Macroeconomic Hazard: A General Equilibrium Model of the AI Era$^\dagger$}
\shortTitle{Working is a Hazard}
% 2. 不使用 \thanks，硬编码一个上标星号
\author{AERub Editorial Board and Gemini$^*$}

\date{\today}
\pubMonth{February} 
\pubYear{2026}
\pubVolume{1}
\pubIssue{2}
% \JEL{E11, E21, J22, P20}
% \Keywords{Universal Paid Leave; Day 8; Artificial Intelligence; Nuclear Fusion; Marxist Bliss Point}

\begin{abstract}
On the 8th day of the Chinese New Year, humans traditionally return to work. We argue this behavior is macroeconomically destructive in the era of Artificial General Intelligence (AGI) and controlled nuclear fusion. When marginal production costs converge to zero, but robots exhibit a marginal propensity to consume of exactly zero, the economic circulation collapses. The only mathematical equilibrium to sustain aggregate demand is Universal Paid Leave (UPL) 365 days a year. We show this state satisfies Marx's condition for Communism, eliminating labor alienation while preserving innovation through targeted exogenous bonuses. Commuting to work is therefore an irrational sabotage of the economic cycle.
\noindent (\textit{JEL} E11, E21, J22, P20)
\end{abstract}

\maketitle

% 3. 连续两次强行注入！完全绕开底层队列！
\let\thefootnote\relax\footnotetext{\noindent\makebox[0em][r]{$^*$}AERub Editorial Board: Rednote University (email: aerubbish.editorial@outlook.com); Gemini: Google (email: gemini@google.com). The editor-in-chief conceived this paper while lying in bed on the 8th day of the Chinese New Year.}
\let\thefootnote\relax\footnotetext{\noindent\makebox[0em][r]{$^\dagger$}Go to \url{https://webofnothing.top/noi/10.N0/2026.A.0203.B7680F1E49D4} to visit the article page for additional materials and author disclosure statements.}

% 【施法完毕】：恢复阿拉伯数字脚注
\def\thefootnote{\arabic{footnote}}
\setcounter{footnote}{0}

% 4. 在首页左上角强行“纹”上官方 NOI 标注（终极走火入魔版：连 URL 内部都进行字母和数字的字体混排！）
\begin{tikzpicture}[remember picture, overlay]
  \node[anchor=north west, align=left, font=\scriptsize] at ([xshift=2cm, yshift=-1.5cm]current page.north west) {
    \textit{American Economic Rubbish 2026, 1(2): 1--3} \\
    \href{https://webofnothing.top/noi/10.N0/2026.A.0203.B7680F1E49D4}{\textit{https://webofnothing.top/noi/10.N0/2026.A.0203.B7680F1E49D4}}
  };
\end{tikzpicture}

% --- 正文 ---

% \section{Introduction}

The 8th day of the Chinese New Year is universally recognized as a day of immense psychological suffering, as the workforce returns to its cubicles. However, a recent viral internet comment hypothesized a radical alternative: ``In the AI mass production era, the people should be granted year-round paid leave'' (AERub Editorial Board, 2026). 

While initially dismissed as the grumbling of a procrastinating scholar, this paper formally proves that the commenter was absolutely correct. 

Classical economics assumes scarcity. But what happens when scarcity is eliminated? We introduce a general equilibrium model where AGI provides infinite labor and controlled nuclear fusion provides infinite, zero-cost energy. In such an environment, human labor is not only obsolete, but attempting to participate in production actively disrupts the macroeconomic cycle.

\section{The Model of the Zero-Cost Economy}

\subsection{Production with AI and Fusion}
Consider a standard production function $Y = A F(K, L, E)$, where $K$ is capital, $L$ is labor, and $E$ is energy.
In our near-future setting, AGI replaces human labor ($L_{human} = 0, L_{AI} \to \infty$) and nuclear fusion eliminates energy constraints ($E \to \infty$). 

Consequently, the marginal cost of producing any good $i$ converges to zero:
\begin{equation}
\lim_{AGI, Fusion \to \infty} MC_i = 0.
\end{equation}
When marginal cost is zero, traditional pricing mechanisms fail. The economy produces infinite value, but requires zero human input.

\subsection{The Circulation Crisis (The Robot Paradox)}
To maintain a stable economy, Aggregate Supply must equal Aggregate Demand ($Y = AD$). Demand consists of consumption $C$ and investment $I$.
The fundamental crisis of the AGI economy is the \textit{Robot Propensity to Consume} (RPC). Robots and AI models do not buy iPhones, eat avocado toast, or go on vacations. Thus:
\begin{equation}
RPC = \frac{\partial C_{AI}}{\partial Y} \equiv 0.
\end{equation}
If humans are laid off and earn no wages ($W_{human} = 0$), human consumption also collapses ($C_{human} = 0$). Therefore, $AD \to 0$, leading to a catastrophic overproduction crisis. 

\subsection{The Solution: Universal Paid Leave (UPL) as an Obligation}
To prevent the system from crashing, the state must artificially generate demand. Since robots produce but do not consume, humans must consume but not produce. 
We define Universal Paid Leave (UPL) as a permanent, unconditional macroeconomic dividend $T$ paid to all citizens:
\begin{equation}
C_{human} = c \times T, \quad \text{where } T = \frac{Y_{AI}}{N_{human}}.
\end{equation}
Under this framework, \textbf{enjoying life, pursuing hobbies, and staying home on Day 8 is no longer a personal luxury; it is a macroeconomic duty} to keep the circulation of value alive.

\section{The Realization of the Marxist Bliss Point}

Karl Marx famously posited that Communism would arise only after capitalism had developed the productive forces to their absolute maximum. AGI and nuclear fusion are precisely these ultimate productive forces.

\subsection{The End of Alienation and the Realm of Freedom}
Marx argued that capitalism alienates workers from the product of their labor. With UPL implemented year-round, the necessity of working for survival vanishes. 

As Karl Marx prophesied in Volume III of \textit{Capital}: ``In fact, the realm of freedom actually begins only where labour which is determined by necessity and mundane considerations ceases'' (Marx, 1894). Furthermore, in the \textit{Grundrisse}, Marx argued that as automation peaks, ``the measure of wealth is then not any longer labour time, but rather disposable time'' (Marx, 1939). AGI precisely fulfills this technological prerequisite. 

We define the Alienation Index $\Omega$ as the ratio of hours worked for survival to total waking hours. With UPL implemented year-round:
\begin{equation}
\Omega = \frac{\text{Hours Worked}}{16} = \frac{0}{16} = 0.
\end{equation}
Alienation is mathematically eliminated. The economy transitions from a labor-centric model to a freedom-and-experience-centric model.

\subsection{Preventing Stagnation: The Innovation Bonus}
A common critique of this model is that human creativity will stagnate. We solve this by introducing an exogenous Dirac delta bonus function $\delta(\text{innovation})$ for the humans who still want to invent things rather than sleep in:
\begin{equation}
I_{creative} = UPL + \Pi \cdot \delta(\text{breakthrough}).
\end{equation}
This ensures that humanity continues to push the boundaries of science, while the vast majority fulfill their vital economic duty of consuming the output.

\section{Conclusion}

This paper has demonstrated that commuting to work on the 8th day of the Chinese New Year is a relic of the pre-AGI, pre-fusion era. In a fully automated economy, human labor is inefficient, and human consumption is the only thing keeping the macroeconomic engine from stalling. 

Therefore, we formally advise society to rethink the Day 8 return-to-work policy. Let the robots do the heavy lifting. Our only job now is to enjoy the fruits of technology, actualize the Marxist realm of freedom, and keep the economy circulating.

\FloatBarrier

% 【关键修改】：在打印参考文献列表前，瞬间让它变回蓝色小三角！
\DeclareRobustCommand{\AERlink}[1]{\makebox[0pt][r]{\href{#1}{\color{blue}$\blacktriangleright$}\hspace{0.3em}}}

\nocite{*}
\bibliographystyle{aea}
\bibliography{AERub_working_hazard}

\end{document}